% ==========================================
% SECTION 9: CONCLUSION AND FUTURE WORK
% ==========================================

\subsection{Summary of Findings}
\label{sub:summary}

This project applied an end-to-end process analytics pipeline to the Road Traffic Fine Management event log, combining process discovery, conformance checking, predictive monitoring, and prescriptive recommendations. The research questions posed in the introduction can now be answered:

\textbf{RQ1 (Process Discovery):} The process is structurally simple — two variants account for over 68\% of all completed cases. The Petri net discovered from the top 10 variants achieves 98.7\% trace fitness on the full log. The primary bottleneck lies in the transition toward credit collection, with mean delays exceeding several hundred days.

\textbf{RQ2 (Operational Behavior):} Batching behavior is clearly visible in the Dotted Chart: administrative activities such as \textit{Send Fine} and \textit{Add penalty} are executed in bulk on specific calendar dates rather than continuously. The Performance Spectrum confirms that early process steps are executed rapidly in synchronized batches, while later steps exhibit long, variable waiting periods driven by citizen response behavior.

\textbf{RQ3 (Predictive Monitoring):} Outcome prediction achieves an AUC-ROC of 0.889 using an LSTM at prefix length $k=5$, though even a simple Logistic Regression at $k=2$ reaches 0.879. The activity sequence alone (CF variant) captures nearly all predictive signal — data-aware features provide only marginal improvement. Remaining time prediction reduces the baseline MAE from 293 to approximately 100 days.

\textbf{RQ4 (Prescriptive Analytics):} SHAP analysis identifies the presence of penalty-related activities and elapsed time as the dominant risk drivers. A threshold-based policy translates predicted probabilities into tiered action recommendations (escalate, remind, or monitor), enabling proactive case management.

\subsection{Limitations}
\label{sub:limitations}

Several limitations constrain the generalizability of this work:

\begin{itemize}
    \item \textbf{Missing resource data:} The event log lacks a usable \texttt{org:resource} attribute, preventing any organizational perspective analysis.
    \item \textbf{Date-only timestamps:} All timestamps are recorded at day granularity, making intra-day performance analysis impossible.
    \item \textbf{Historical data:} The log covers 2000–2013. Process behavior, legal frameworks, and administrative systems may have changed since then.
    \item \textbf{Cost parameters:} The prescriptive recommendations rely on illustrative cost assumptions that would require calibration with real operational data.
    \item \textbf{Model saturation:} The near-identical performance across model families suggests that the prediction problem is largely solved by simple features. More complex architectures do not yield meaningful improvements on this dataset.
\end{itemize}

\subsection{Future Work}
\label{sub:future_work}

Possible extensions include: applying Transformer-based sequence models to evaluate whether attention mechanisms capture long-range dependencies better than LSTMs; implementing online process mining for real-time case monitoring as events arrive; integrating actual cost data from the administration to calibrate the prescriptive policy; and extending the approach to other administrative domains where similar fine-to-collection escalation processes exist.